\section*{الفصل \ref{ch:b3:membrane-traffic} --- حركة الأغشية وفرز البروتينات}

\begin{solution}{exo:b3:membrane-traffic:1}
شوط أميني الطرف كاره للماء: \omterm{def:b3:membrane-traffic:signals}{ببتيد إشارة}، إلى
الشبكة الهيولية الباطنة (ومنها افتراضًا إلى السطح). وPKKKRKV:
إشارة توطين نووي، إلى النواة عبر المسام.
ولولب أميني الطرف برمائي: \omterm{def:b3:membrane-traffic:signals}{تسلسل أمامي} متقدري، إلى
حشوة المتقدرة. وKDEL: استرجاع من غولجي رجوعًا إلى
الشبكة. ومانوز-6-فوسفات: من غولجي البعيد إلى \omterm{def:b3:membrane-traffic:endocytosis}{الجسيم الداخلي}
\omterm{def:b3:membrane-traffic:endocytosis}{والجسيم الحال}.
\end{solution}

\begin{solution}{exo:b3:membrane-traffic:2}
(1) من دون أغشية كانت السلسلة المترجمة أطول من البروتين
المفرَز بنحو عشرين بقية: أي أن قطعة تُزال أثناء
الإفراز. (2) وبحضور الأغشية أثناء الترجمة كان للناتج
الحجم الناضج وكان محميًا من البروتياز: أي أنه دخل
الحويصلات وفقد القطعة داخلها. (3) والأغشية المضافة
بعد الترجمة لم تفعل أيًا من ذلك: فالدخول مقترن بالتركيب
--- أي مصاحب للترجمة --- والقطعة، وهي \omterm{def:b3:membrane-traffic:signals}{ببتيد الإشارة}، لا تعمل إلا
على سلسلة وليدة.
\end{solution}

\begin{solution}{exo:b3:membrane-traffic:3}
الشبكة إلى غولجي: \omterm{def:b3:membrane-traffic:coats}{COPII}، إلى الأمام. وغولجي إلى الشبكة: \omterm{def:b3:membrane-traffic:coats}{COPI}،
إلى الوراء. وغولجي البعيد إلى \omterm{def:b3:membrane-traffic:endocytosis}{الجسيم الداخلي}: \omterm{def:b3:membrane-traffic:coats}{كلاذرين} بمهايئاته،
إلى الأمام نحو \omterm{def:b3:membrane-traffic:endocytosis}{الجسيم الحال}. والغشاء الهيولي إلى \omterm{def:b3:membrane-traffic:endocytosis}{الجسيم الداخلي}: \omterm{def:b3:membrane-traffic:coats}{كلاذرين}،
إلى الداخل (التقام).
\end{solution}

\begin{solution}{exo:b3:membrane-traffic:4}
لا تُضاف السكريات إلا داخل لمعة الشبكة وغولجي. و
الوجه اللمعي لكل حويصلة يبقى لمعيًا عبر التبرعم والاندماج،
وحين تندمج الحويصلة بالغشاء الهيولي يصير وجهها اللمعي
خارج الخلية؛ فالاتجاه لا ينقلب أبدًا، ومن ثم
ينتهي ما غُلكز في الداخل إلى الخارج.
\end{solution}

\begin{solution}{exo:b3:membrane-traffic:5}
$t^{*} = \ln(0.05/0.2)/(0.05 - 0.2) = \ln 4/0.15 = \qty{9.2}{min}$؛
و$B(t^{*}) = (0.2/(-0.15))(e^{-1.85} - e^{-0.46}) = 0.63$. وعند \qty{60}{min}
يكون $A \approx 0$، و$B = 1.33\,(e^{-3} - e^{-12}) = 0.066$، ومنه $C = 0.93$.
\end{solution}

\begin{solution}{exo:b3:membrane-traffic:6}
\omterm{def:b3:membrane-traffic:signals}{ببتيد الإشارة} يضع النهاية الأمينية في اللمعة؛ وكل
شوط كاره للماء تسلسلُ إيقاف نقل أو بدء نقل بالتناوب،
فتكون اللوالب العابرة للغشاء أربعة (نحو البقايا 30--52،
و90--112، و150--172، و210--232). والقطع: النهاية الأمينية لمعية؛
و52--90 عصارية؛ و112--150 لمعية؛ و172--210 عصارية؛ والنهاية
الكربوكسيلية لمعية. والغلكزة على القطع اللمعية وحدها --- أي
النهاية الأمينية، و112--150 والنهاية الكربوكسيلية --- وهي التي تصير
خارج خلوية عند السطح.
\end{solution}

\begin{solution}{exo:b3:membrane-traffic:7}
$J_{\max} = \num{30000}\times 0.25\times 0.05/0.30 = 1250$ في الدقيقة؛
ونسبة السطح $k_{r}/(k_{i} + k_{r}) = 0.05/0.30 = 0.17$. ومضاعفة $R$
تضاعف $J$ بالضبط؛ ومعدل العودة هو العنق (فمعظم المستقبلات
في الداخل)، ومضاعفة $k_{r}$ تعطي $\num{30000}\times 0.25
\times 0.1/0.35 \approx 2140$، أي بعامل $1.7$.
\end{solution}

\begin{solution}{exo:b3:membrane-traffic:8}
(أ) بلا مستقبل: تتبع الحلمهات الموسومة الطريق الافتراضي
وتُفرَز؛ فتخلو الجسيمات الحالة من الإنزيمات وتمتلئ بالمواد
غير المهضومة --- أي داء خزن. (ب) وبلا ناقل فوسفاتي: النمط
الظاهري نفسه بآفة أخرى (داء الخلايا I)، فالإنزيمات
غير موسومة وتُفرَز. (ج) وبجسيمات داخلية معادَلة: لا يستطيع المستقبل
أن يطلق حمولته، فلا يُعاد تدويره وتُوصَّل الحلمهات
خطأً أو تُفرَز؛ وتتوقف مستقبلات LDL وغيرها عن الدوران أيضًا.
\end{solution}

\begin{solution}{exo:b3:membrane-traffic:9}
التيروزين ينتمي إلى \omterm{def:b3:bioinformatics:motif}{نميطة} الذيل التي يتعرفها مهايئ
\omterm{def:b3:membrane-traffic:coats}{الكلاذرين}؛ ومن دونه لا يُجمع المستقبل في الحفر المغطاة.
فتربط المستقبلات LDL طبيعيًا لكنها تنتشر انتشارًا منتظمًا على
السطح، مستبعدةً من الحفر، وتبقى الرابطة في الخارج --- أي
داء فرز لا داء ارتباط.
\end{solution}

\begin{solution}{exo:b3:membrane-traffic:10}
عند $k_{1} = k_{2} = k$، جرّب $B = kt\,e^{-kt}$، فيكون $\mathrm{d}B/\mathrm{d}t
= k e^{-kt} - k^{2}t e^{-kt} = kA - kB$، كما هو مطلوب، مع $B(0) = 0$.
والذروة حيث $1 - kt = 0$: أي $t^{*} = 1/k$، و$B = e^{-1} = 0.37$. وعمومًا
$C(t) = 1 - (k_{2}e^{-k_{1}t} - k_{1}e^{-k_{2}t})/(k_{2} - k_{1})$،
وهي متناظرة عند تبادل $k_{1}$ وَ$k_{2}$: فمنحني الحبيبات
وحده لا يستطيع أن يقول أي المقصورتين هي السريعة.
\end{solution}

\begin{solution}{exo:b3:membrane-traffic:11}
مساحة كل حويصلة $4\pi(50\,\text{nm})^{2} = \qty{0.031}{\micro
m^2}$؛ وألف في الدقيقة تعطي \qty{31}{\micro m^2/min}، فيلزم للسطح
كله $3000/31 \approx \qty{95}{min}$. ولا تنكمش الخلية
لأن الغشاء يعود بالإخراج الخلوي (حويصلات إعادة التدوير) بالمعدل
نفسه --- أي توازن تدفقات. ويسع كل حويصلة $5\times 10^{-19}$ لتر؛ وألف
في الدقيقة على مدى ساعة تعطي $3\times 10^{-14}$ لتر، أي نحو
\qty{1.5}{\%} من حجم خلية قدره \qty{2}{pL} في الساعة.
\end{solution}

\begin{solution}{exo:b3:membrane-traffic:12}
ذيفان الوشيقية يعمل في النهاية الحركية: فلا يُطلق أستيل كولين،
ولا تتلقى العضلة أمرًا فترتخي --- أي شلل
رخو. أما ذيفان الكزاز فيُحمل صعودًا في المحور الحركي وعبره إلى
العصبونات البينية المثبطة التي تكبح العصبونات الحركية عادةً؛ فإذا
قُطع بروتين \omterm{def:b3:membrane-traffic:fusion}{SNARE} فيها لم يُطلق غليسين ولا غابا، فتطلق العصبونات الحركية
بلا كابح وتتقلص كل عضلة --- أي شلل تشنجي. وبضعة
نانوغرامات من ذيفان الوشيقية محقونة في عضلة مفرطة النشاط
تُسكتها موضعيًا شهورًا (حتى يُصنع \omterm{def:b3:membrane-traffic:fusion}{SNARE} جديد وتتعافى النهاية):
وهو علاج لخلل التوتر والتشنج والحول والتجاعيد.
\end{solution}

\begin{solution}{pb:b3:membrane-traffic:1}
\textbf{1.} $10^{7}\times \qty{25000}{Da}\times 1.66\times 10^{-24}$ غرام
$= \qty{0.42}{pg}$ في الدقيقة، أي \qty{600}{pg} في اليوم --- أي ضعف محتواها
من البروتين.
\textbf{2.} حجم الحويصلة $\tfrac{4}{3}\pi\,35^{3} = 1.8\times 10^{5}$
نانومتر$^{3}$، و\qty{30}{\%} منه $5.4\times 10^{4}$؛ وجزيء الإنزيم الواحد
$\tfrac{4}{3}\pi\,2^{3} = 34$ نانومتر$^{3}$: أي نحو \num{1600} جزيء في
الحويصلة؛ و$10^{7}/1600 \approx 6200$ حويصلة في الدقيقة.
\textbf{3.} المساحة $4\pi\,35^{2} = \qty{0.0154}{\micro m^2}$ لكل واحدة، أي \qty{95}{\micro
m^2} في الدقيقة: فشبكةُ \qty{6000}{\micro m^2} في نحو
ساعة لو لم يعد شيء.
\textbf{4.} $\qty{95}{\micro m^2} = 9.5\times 10^{7}$ نانومتر$^{2}$، وبطبقتين
عند \qty{0.7}{nm^2}: أي $2.7\times 10^{8}$ دهن في الدقيقة. والخلية
تعيد الغشاء بحويصلات \omterm{def:b3:membrane-traffic:coats}{COPI} وبإعادة التدوير، وتركّب دهنًا
لتعويض ما يغادر نهائيًا في الحبيبات.
\textbf{5.} حجم الحبيبة $5.2\times 10^{8}$ نانومتر$^{3}$، و\qty{30}{\%} منه
$1.6\times 10^{8}$: أي $4.7\times 10^{6}$ جزيء في الحبيبة؛ و$6\times
10^{8}$ جزيء في الساعة تملأ نحو $130$ حبيبة.
\textbf{6.} كل حبيبة تضيف $\pi d^{2} = \qty{3.1}{\micro m^2}$؛ و$300$
تضيف \qty{940}{\micro m^2} إلى \qty{50}{\micro m^2}: أي عشرين ضعفًا. و
يجب أن يُسترجع الغشاء بالتقام تعويضي في دقائق.
\textbf{7.} $t^{*} = \ln 2/0.05 = \qty{13.9}{min}$، و$B = 0.50$.
\textbf{8.} $A = e^{-3} = 0.05$؛ و$B = 2(e^{-1.5} - e^{-3}) = 0.35$؛
و$C = 0.60$.
\textbf{9.} \qty{10}{min} و\qty{20}{min}؛ والمجموع الوسطي \qty{30}{min}.
\textbf{10.} $t^{*} = \ln(0.1)/(-0.09) = \qty{25.6}{min}$؛ و$B(t^{*}) =
(0.1/(-0.09))(e^{-2.56} - e^{-0.256}) = 0.77$: فينتفخ غولجي،
محتقنًا بثلاثة أرباع الحمولة.
\textbf{11.} عند الذروة $k_{1}e^{-k_{1}t^{*}} = k_{2}e^{-k_{2}t^{*}}$؛
وبالتعويض في $B$ نجد $B_{\max} = (k_{1}/k_{2})^{\,k_{2}/(k_{2} -
k_{1})}$: فالأسّ هو $k_{2}/(k_{2} - k_{1})$. ومن أجل $k_{1} > k_{2}$ يكون
الأسّ سالبًا والأساس فوق $1$؛ ومن أجل $k_{1} < k_{2}$ يكون
الأساس دون $1$ والأسّ موجبًا --- وفي الحالتين تكون القيمة
دون $1$ (وللتحقق: $2^{-1} = 0.5$ هنا).
\textbf{12.} من أجل $k_{1} \gg k_{2}$، تؤول $B(t) \to e^{-k_{2}t}$: فالواسم
يكون في غولجي فورًا ويغادره بمعدل $k_{2}$ --- فخطوة الشبكة
غير مرئية ومنحني غولجي تلاشٍ بسيط.
\textbf{13.} $\theta = 2/(2+2) = 0.5$؛ و$J = \num{50000}\times 0.2\times
0.1\times 0.5/(0.1 + 0.1) = 2500$ جسيم في الدقيقة، أي \num{150000} في
الساعة.
\textbf{14.} $1.5\times 10^{5}\times 1500 = 2.3\times 10^{8}$ جزيء
كوليسترول في الساعة؛ و$5\times 10^{9}/2.3\times 10^{8} \approx 22$ ساعة.
\textbf{15.} نسبة السطح $k_{r}/(k_{r} + k_{i}\theta) = 0.5$؛ والجولة
الواحدة تستغرق $1/(k_{i}\theta) + 1/k_{r} = 10 + 10 = \qty{20}{min}$:
أي نحو $60$ جولة في عشرين ساعة.
\textbf{16.} $J = 1250$ في الدقيقة. ودون الإشباع يكون $J \propto R\,[L]$،
فبنصف المستقبلات يجب أن يتضاعف LDL البلازما للتصفية نفسها.
\textbf{17.} تعود $R$ إلى \num{50000}: فتُستعاد التصفية، ويهبط LDL البلازما
عائدًا نحو السوي --- وهو أثر الستاتين.
\textbf{18.} بلا مستقبلات لا يُصفّى LDL إلا بطرق
غير نوعية بطيئة، فيدور أيامًا بدل ساعات، ويتأكسد،
وتلتقطه المستقبلات الكانسة في البلعميات،
فتصير الخلايا الرغوية في اللويحات العصيدية؛ كما أن الستاتينات،
التي تعمل بتحريض المستقبلات، لا تجد ما تحرضه.
\textbf{19.} $V = \tfrac{4}{3}\pi(0.25)^{3} = \qty{0.065}{\micro m^3} =
6.5\times 10^{-17}$ لتر. وعند pH $7.2$: $6.3\times 10^{-8}\times 6.5\times
10^{-17} = 4\times 10^{-24}$ مول --- أي نحو $2$ بروتون حر. وعند pH $4.7$:
$2\times 10^{-5}\times 6.5\times 10^{-17} = 1.3\times 10^{-21}$ مول،
أي نحو $780$.
\textbf{20.} $50\times 780 \approx \num{39000}$ بروتون؛ و$50$ مضخة عند
$200$ في الثانية تنقل \num{10000} في الثانية: أي نحو \qty{4}{s}.
\textbf{21.} ضخّ الشحنة الموجبة إلى الداخل يجعل اللمعة موجبة، و
بضعة آلاف من الشحنات توقف المضخة؛ فيدخل الكلور عبر
قناة (أو تخرج كاتيونات) لتعادل الشحنة.
\textbf{22.} الإنزيمات الخاملة عند pH المتعادل لا تضر إذا رشح \omterm{def:b3:membrane-traffic:endocytosis}{جسيم حال}
أو إذا فُرزت خطأً إلى العصارة أو إلى السطح: فاشتراط الحمض
يحصر الهضم في المقصورة المبنية له.
\textbf{23.} لا بد أن يحمل كلاهما مستشعر pH --- هستيدينات
تغيّر بَرْتَنَتُها قرب pH~6 سطحَ الارتباط: فمستقبل LDL
يطوي نطاقًا فوق موقع ارتباطه نفسه عند البرتنة، و
مستقبل مانوز-6-فوسفات يفقد ألفته كذلك.
\textbf{24.} $10^{7.2 - 4.7} = 10^{2.5} \approx 300$ ضعف تركيزًا؛
والقاعدة المحتجَزة تستهلك البروتونات، فترفع pH \omterm{def:b3:membrane-traffic:endocytosis}{الجسيم الحال}، و
تعطّل الحلمهات --- وهكذا يقتل الكلوروكين طفيليات الملاريا
في فجوتها الغذائية وهكذا يحجب \omterm{def:b3:membrane-traffic:endocytosis}{الالتهام الذاتي} في التجارب.
\textbf{25.} ذروة غولجي عند \qty{14}{min}؛ ونحو \num{6200} حويصلة تغادر
الشبكة في الدقيقة؛ و\num{150000} جسيم LDL ملتقَط في
الساعة.
\end{solution}
